\documentclass[10pt]{article}
\usepackage[T1]{fontenc}
\usepackage[utf8]{inputenc}
\usepackage[french]{babel}
\usepackage{amsmath,amssymb,mathrsfs,mathpazo}
\usepackage{geometry}
\geometry{a4paper,margin=2cm}
\usepackage{graphicx,array,tabularx,tabularray,colortbl,multirow,multicol,booktabs,enumitem,pifont,ulem,fancybox,diagbox,cancel}
\usepackage{pstricks,pstricks-add,pst-plot,pst-tree,pst-node,pst-text,pst-func,pst-math,pst-bspline,pst-3dplot,pst-eucl,pst-solides3d}
\let\pscalculate\undefined
\usepackage{pst-bezier}
\usepackage{tikz,pgfplots,tkz-tab}
\usetikzlibrary{arrows,arrows.meta,patterns,calc,positioning,trees,decorations.markings,shapes.geometric,shapes.misc,shapes.arrows,intersections,decorations.pathmorphing,decorations.pathreplacing,backgrounds,shadows,babel}
\pgfplotsset{compat=1.18}
\usepgfplotslibrary{fillbetween}
\usepackage[np]{numprint}
\usepackage{listings,makecell,pas-tableur,twoopt}
\let\barre\undefined
\usepackage[pstricks]{tablvar}
\definecolor{darkgreen}{RGB}{0,100,0}
\providecommand{\red}{\color{red}}
\providecommand{\blue}{\color{blue}}
\providecommand{\green}{\color{green}}
\providecommand{\black}{\color{black}}
\providecommand{\white}{\color{white}}
\providecommand{\gray}{\color{gray}}
\providecommand{\cyan}{\color{cyan}}
\providecommand{\magenta}{\color{magenta}}
\providecommand{\yellow}{\color{yellow}}
\providecommand{\violet}{\color{violet}}
\providecommand{\orange}{\color{orange}}
\providecommand{\euro}{\texteuro}
\providecommand{\EUR}{\texteuro}
\providecommand{\ieme}{\textsuperscript{e}}
\providecommand{\degres}{\ensuremath{^\circ}}
\providecommand{\degree}{\ensuremath{^\circ}}
\providecommand{\vv}[1]{\overrightarrow{#1}}
\providecommand{\Coord}[1]{#1}
\providecommand{\decosix}{}
\providecommand{\starredbullet}{\ensuremath{\bullet}}
\pagestyle{empty}
\setlength{\parindent}{0pt}
\usepackage{adjustbox}\DeclareUnicodeCharacter{2236}{:}\DeclareUnicodeCharacter{2190}{\ensuremath{\leftarrow}}
\providecommand{\R}{\mathbb{R}}
\providecommand{\N}{\mathbb{N}}
\providecommand{\Z}{\mathbb{Z}}
\providecommand{\Q}{\mathbb{Q}}
\providecommand{\D}{\mathbb{D}}
\providecommand{\C}{\mathbb{C}}
\providecommand{\e}{{\mathrm{e}}}
\providecommand{\vect}[1]{\overrightarrow{#1}}
\providecommand{\vectt}[1]{\overrightarrow{\mathrm{#1}}}
\providecommand{\vvt}[1]{\overrightarrow{\mathrm{#1}}}
\providecommand{\np}[2][]{{#2~#1}}
\providecommand{\vv}[1]{{\overrightarrow{#1}}}
\renewcommand{\d}{{\mathrm{d}}}
\providecommand{\degres}{{^\circ}}
\providecommand{\degre}{{^\circ}}
\providecommand{\blue}{\color{blue}}
\providecommand{\red}{\color{red}}
\providecommand{\enskip}{\quad}
\providecommand{\llbracket}{\lbrack\!\lbrack}
\providecommand{\rrbracket}{\rbrack\!\rbrack}
\providecommand{\euro}{{\text{€}}}
\providecommand{\no}{{\text{n°}}}
\providecommand{\interval}[2]{{[#1;#2]}}
\providecommand{\intervFO}[2]{{[#1;#2[}}
\providecommand{\dotfill}{\ldots}

\providecommand{\ds}{\displaystyle}
\providecommand{\textding}[1]{#1}
\providecommand{\Oij}{\(\mathrm{O};\vec{\imath},\vec{\jmath}\)}
\providecommand{\Oijk}{\(\mathrm{O};\vec{\imath},\vec{\jmath},\vec{k}\)}
\providecommand{\Ouv}{\(\mathrm{O};\vec{u},\vec{v}\)}
\providecommand{\intervFF}[2]{\ensuremath{[#1;#2]}}
\renewcommand{\euro}{\ifmmode\text{\texteuro}\else\texteuro\fi}
\DeclareUnicodeCharacter{00B0}{\ifmmode^{\circ}\else\textdegree\fi}
\DeclareUnicodeCharacter{20AC}{\ifmmode\text{\texteuro}\else\texteuro\fi}
% Énoncés : Éducation nationale. Transcription et corrigés : APMEP, Denis Vergès et contributeurs. Découpage et classement : CoopMaths, Sébastien Lozano, Éric Elter et contributeurs. Adaptation de la présentation : Maths Explorer.
% CC BY-SA : https://coopmaths.fr/www/about/
% https://coopmaths.fr/alea/static/bac/2026/tex/bac_2026_06_sujet2_polynesie_4_cor.tex
\begin{document}
\noindent\textit{APMEP / CoopMaths. CC BY-SA. Présentation adaptée par Maths Explorer.}\par\medskip
\begin{enumerate}
		\item Soit $g(x) = \e^{\left(x^2\right)}$

		Avec $u(x) = x^2$ puis $u'(x) = 2x$, on a

		$\left(\e^{\left(x^2\right)}\right)' = \left(\e^{u(x))}\right)' = \e^{u(x)} \times u'(x) = \e^{\left(x^2\right)} \times 2x$.

		Donc $f'(x) = \e^{\left(x^2\right)} + x \times \e^{\left(x^2\right)} \times 2x$ ;

		$f'(x) = \e^{\left(x^2\right)}\left(1 + 2x^2\right)$.

		Puis de même $f''(x) = 4x\e^{\left(x^2\right)} + \left(1 + 2x^2\right) \times 2x \e^{\left(x^2\right)} = 2x \left(2 + 1 + 2x^2 \right)\e^{\left(x^2\right)}$.

		Finalement $f''(x) = 2x\left(3 + 2x^2 \right)\e^{\left(x^2\right)}$

		Dans ce produit pour $x \in ]- \infty ; 0[$, on a :

		\begin{itemize}
			\item $3 + 2x^2 > 0$ ;
			\item $\e^{\left(x^2\right)} > 0$
			\item $2x < 0$
		\end{itemize}

		Conclusion sur $]-\infty ; 0]$, \: $f''(x) < 0$ : la fonction $f$ est concave sur cet intervalle.

		\hfill\textbf{Affirmation 1 vraie}

		\item \begin{itemize}
			\item Dans le groupe 1 on peut former $\ds\binom{24}{8}$ ensembles de 8 rapaces ;

			\item Dans le groupe 2 on peut former $\ds\binom{16}{8}$ ensembles de 8 rapaces ;

			\item Dans le groupe 3 on peut former $\ds\binom{8}{8}$ ensemble de 8 rapaces ;
		\end{itemize}

		Il y a donc $\ds\binom{24}{8} \times \binom{16}{8} \times \binom{8}{8}$ façons différentes de répartir 8 rapaces dans ces 3 groupes.

		Or\quad $\ds \binom{24}{8} \times \binom{16}{8} \times \binom{8}{8} = \dfrac{24!}{8! \times (24- 8)!} \times \dfrac{16!}{8!\times (16-8)!} \times 1
		= \dfrac{24!}{8! \times 16!} \times \dfrac{16!}{8! \times 8!  }
		= \dfrac{24!}{\left (8!\right )^3}$.

		\hfill\textbf{Affirmation 2 vraie}

		\item $f$ est solution de l'équation différentielle donc quel que soit $x \in \R$,

		$f'(x) = - f(x) + 1$ \quad (1) .

		D'autre part $g(x) = f(x) - 1$ montre que puisque $f$ est dérivable sur \, $g$ l'est aussi avec : $g'(x) = f'(x)$.

		L'égalité (1) devient alors

		$g'(x) 	= - f(x) - 1 = - (f(x) - 1) = - g(x)$

		Conclusion : quel que soit $x \in \R, \: g'(x)= - g(x)$.

		\hfill\textbf{Affirmation 3 vraie}

		\item Quel que soit $n \in \N^*, \: - 1 \leqslant \sin (n) \leqslant 1$.

		En multipliant par $\dfrac1n > 0$, on a
		$- \dfrac{1}{n} \leqslant \dfrac{\sin(n)}{n} \leqslant \dfrac{1}{n}$ .

		Comme $\displaystyle\lim_{n \to + \infty} -\dfrac{1}{n} = \displaystyle\lim_{n \to + \infty} \dfrac{1}{n} = 0$, on a donc d'après le théorème des gendarmes:
		\[\displaystyle\lim_{n \to + \infty}\dfrac{\sin(n)}{n} = 0.\]

		La suite converge vers zéro.

		\hfill\textbf{Affirmation 4 fausse}

		\item ~
		\begin{center}
			
			
			
			
			
			
			
			
			
			
			
			
			
			
			
			
			
			
			\psset{unit=1cm,arrowsize=2pt 3}
			\begin{pspicture*}(-1.1,-1.5)(5.5,5.5)
				\psaxes[linewidth=1.25pt,labelFontSize=\scriptstyle]{->}(0,0)(-1,-1.5)(5.5,5.5)
				\psplot[plotpoints=2000,linewidth=1.25pt,linecolor=red]{-0.5}{4.5}{4 x mul 1 add x dup mul sub}
				\psplot[plotpoints=500,linewidth=1.25pt,linecolor=blue]{-1}{5}{x 1 add}
				\pscustom[fillstyle=solid,fillcolor=lightgray]{\psplot[plotpoints=2000,linewidth=1.25pt,linecolor=red]{0}{3}{4 x mul 1 add x dup mul sub}\psline(3,4)(0,1)}
				
				\psgrid[gridlabels=0pt,subgriddiv=2,gridwidth=0.2pt,subgridwidth=0.1pt]
				\pspolygon[fillstyle=vlines](0,0)(0,1)(3,4)(3,0)
				\uput[ul](0,0){O}\uput[ul](0,1){I}\uput[r](3,4){A}\uput[dr](3,0){B}
			\end{pspicture*}
			
		\end{center}

		L'aire du domaine grisé est égale à la différence de deux aires :

		\begin{itemize}
		\item l'aire de la surface limitée par l'axe des abscisses, les droites d'équation $x = 0$ et $x = 4$ et la représentation graphique de la fonction définie sur $\R$ par \quad $f : x \longmapsto - x^2 + 4x + 1$.

		La fonction $f$ étant manifestement positive sur l'intervalle $[0 ; 3]$, cette aire est l'intégrale (en  u. a.) $\displaystyle\int_0^3 f(x) \:\text{d}x = \displaystyle\int_0^3 - x^2 + 4x + 1)\:\text{d}x = \left[-\dfrac{x^3}{3} + 2x^2 + x \right]_0^3 = -\dfrac{3^3}{3} + 2\times 3^2 + 3 - \left(-\dfrac{0^3}{3} + 2\times 0^2 + 0 \right) = - \dfrac{27}{3} + 18 + 3 = 22 - 9 = 12$ \: (u.a.) \:\textbf{et}

		\item l'aire du trapèze rectangle OIAB soit $\dfrac{\text{OI} + \text{AB}}{2}\times \text{OB}$;
		\end{itemize}

		On a de façon évidente : OI $ = 1$, \: AB $ = 4$ et OB $ =  3$, donc

		$\mathcal{A}(\text{OIAB}) = \dfrac{1 + 4} {2} \times 3 = \dfrac{15}{2}$\: (u.a)

		Si $\mathcal{A}(G)$ est l'aire de la surface grisée alors

		$\mathcal{A}(G) = 12 - \dfrac{15}{2} = \dfrac{24 - 15}{2} = \dfrac{9}{2} = 4,5$\:(u.a.).

		\hfill\textbf{Affirmation 5 fausse}

		\medskip

		\emph{Remarque} : on peut aussi dire que $\mathcal{A}(G)$  est la différence de l'intégrale de la fonction $f(x) = - x^2 + 4x + 1$ et de l'intégrale de la fonction $g(x) = x + 1$ sur l'intervalle [0 ; 3].

		$\mathcal{A}(G) = \displaystyle\int_0^3 \left(- x^2 + 4x + 1\right)\:\text{d}x - \displaystyle\int_0^3 (x + 1)\:\text{d}x$, soit par linéarité de l'intégrale :

		$\mathcal{A}(G) = \displaystyle\int_0^3 \left(- x^2 + 4x + 1 - (x + 1) \right)\:\text{d}x =
		\displaystyle\int_0^3 \left(- x^2 + 4x + 1 - x - 1) \right)\:\text{d}x = \displaystyle\int_0^3 \left(- x^2 + 3x  \right)\:\text{d}x =
		\displaystyle\int_0^3 x(3 - x) :\text{d}x$.

		Il faut alors considérer cette fonction trinôme $x \longmapsto x(x - 3)$ qui est négative sauf entre ses racines 0 et 3 sui est justement l'intervalle d'intégration.

		Il faut alors rappeler que cette fonction polynôme trinôme est continue sur $\R$ donc en particulier sur [0 ; 3], donc intégrale sur cet intervalle et affirmer que

		$\mathcal{A}(G) = \displaystyle\int_0^3 \left(- x^2 + 3x \right)\:\text{d}x = \left[- \dfrac{x^3}{3} + 3\frac{x^2}{2} \right]_0^3 = -\dfrac{3^3}{3} + 3\times \dfrac{3^2}{2} - \dfrac{0^3}{3} - 3\times \dfrac{0^2}{2}  = -9 + \dfrac{27}{2} = \dfrac{-18 + 27}{2} = \dfrac92 \ne \dfrac{13}{2}$.
\end{enumerate}

\end{document}
